\documentclass[11pt]{article}
\usepackage[letterpaper,margin=0.82in]{geometry}
\usepackage{amsmath,amssymb,amsthm,mathtools}
\usepackage{graphicx}
\usepackage[hidelinks]{hyperref}
\usepackage{microtype}
\usepackage[T1]{fontenc}

\newtheorem{theorem}{Theorem}
\newtheorem{lemma}{Lemma}
\newcommand{\A}{\mathbb A}
\newcommand{\one}{\mathbf 1}
\newcommand{\HS}{\operatorname{HS}}
\newcommand{\ocl}{\operatorname{ocl}}
\newcommand{\AP}{\operatorname{AP}}
\newcommand{\ran}{\operatorname{ran}}
\newcommand{\mix}{\operatorname{mix}}

\title{Compactness implies discrete spectrum for arbitrary\\
Kaplansky--Hilbert dynamical systems}
\author{Candidate full solution to Remark 3.8 of arXiv:2505.14132}
\date{13 August 2026}

\begin{document}
\maketitle

\paragraph{Status.}
\textbf{Candidate full solution; likely valid, pending expert review.}
The result answers affirmatively the implication left open in Remark 3.8 of
Haase--Kreidler~\cite{HK}: for an arbitrary KH-dynamical system, order-density
of the conditionally almost-periodic vectors implies discrete spectrum.  In
fact, every such vector already belongs to the discrete-spectrum part.

\section{Source question}

Let $(\A,S;E,T)$ be a KH-dynamical $G$-system: $\A$ is a Stone algebra,
$E$ is a Kaplansky--Hilbert module over $\A$, $S:G\to\operatorname{Aut}(\A)$,
and each $T_t$ is an $S_t$-covariant unitary.  Write
\[
 E_{\rm ds}=\ocl\!\bigcup\{M\subseteq E:M\text{ is a finite-rank,
 $T$-invariant KH-submodule}\}.
\]
A vector $x\in E$ is \emph{conditionally almost periodic} if, for every
$\varepsilon>0$, there is a finite $F\subseteq E$ such that
\begin{equation}\label{eq:ap}
 \inf_{y\in F}|T_t x-y|\leq \varepsilon\one\qquad(t\in G).
\end{equation}
The source asks whether order-density of these vectors forces
$E=E_{\rm ds}$~\cite[Remark 3.8]{HK}.  The source proves this for KH-systems
coming from measure-preserving extensions, using an ambient global
$L^2$-space, but explicitly leaves the abstract KH case open.

\begin{center}
\includegraphics[width=0.98\textwidth]{figures/open_problem_crop.png}
\end{center}
\noindent\emph{Source evidence:} arXiv:2505.14132v3, PDF p.~14, Remark 3.8.

\section{Proof intuition}

Condition~\eqref{eq:ap} supplies a finite-rank module whose orthogonal
projection $P$ almost contains the whole orbit of $x$.  Conjugating $P$ back
along the dynamics produces finite-rank projections $Q_t$ which all almost
fix the \emph{same} vector $x$.  We take the order-closed cyclic convex hull
of the $Q_t$ inside the KH-module $\HS(E)$.  A lattice-valued version of the
minimum-norm argument for Hilbert spaces gives a fixed point in this hull.
It is an invariant Hilbert--Schmidt operator $A_\varepsilon$, and the hull
construction preserves
$|x-A_\varepsilon x|\leq\varepsilon\one$.  The equivariant spectral theorem
puts the range of every invariant Hilbert--Schmidt operator in $E_{\rm ds}$.
Letting $\varepsilon\downarrow0$ proves $x\in E_{\rm ds}$.

The cyclic (mix-closed) hull is the crucial feature: it lets one select the
smaller of two $\A$-valued norms on complementary idempotent pieces, which is
what makes the minimum-norm element exist.

\section{A cyclic convex fixed-point lemma}

For a bounded subset $C$ of a KH-module, call $C$ \emph{mix-closed} if
every mixing $\sum_i p_i z_i$, with $(p_i)$ a partition of unity by
idempotents and $z_i\in C$, belongs to $C$.

\begin{lemma}[cyclic convex fixed point]\label{lem:fixed}
Let $(\A,R;H,U)$ be a KH-dynamical $G$-system.  Every nonempty bounded,
order-closed, convex, mix-closed, $U$-invariant subset $C\subseteq H$
contains a fixed point of $U$.
\end{lemma}

\begin{proof}
Set
\[
 h:=\inf_{z\in C}|z|^2\in\A_+.
\]
The set of squared norms is downward directed.  Indeed, for $u,v\in C$,
the Stone property supplies an idempotent $p\in\A$ which selects the two
order regions so that
\[
 p|u|^2+p^c|v|^2=|u|^2\wedge|v|^2.
\]
(Under $\A=C(\Omega)$, take the clopen closure of the region
$|u|^2<|v|^2$; on its boundary the functions agree.)  Since $C$ is
mix-closed, $w=pu+p^cv$ belongs to $C$, and
$|w|^2=|u|^2\wedge|v|^2$.

For every nonempty finite $K\subseteq C$, choose $z_K\in C$ with
\[
 r_K:=|z_K|^2=\bigwedge_{z\in K}|z|^2.
\]
Ordered by inclusion, $r_K\downarrow h$.  Convexity puts
$(z_K+z_L)/2$ in $C$, so the parallelogram identity gives
\begin{align*}
 |z_K-z_L|^2
 &=2r_K+2r_L-4\left|\frac{z_K+z_L}{2}\right|^2\\
 &\leq 2(r_K-h)+2(r_L-h).
\end{align*}
For $K,L\supseteq K_0$, the last expression is at most
$4(r_{K_0}-h)$, which decreases to zero.  Hence $(z_K)$ is
order-Cauchy.  Order completeness and order-closedness give
$z_K\to z_0\in C$ in order, and order continuity of the modulus yields
$|z_0|^2=h$.

This minimizer is unique: if $|u|^2=|v|^2=h$, applying the same identity
to their midpoint gives $|u-v|^2\leq0$.  Finally, $U_tC=C$ and covariance
imply
\[
 R_t h
 =\inf_{z\in C}R_t|z|^2
 =\inf_{z\in C}|U_tz|^2=h.
\]
Thus $U_tz_0$ is another minimizer, and uniqueness gives
$U_tz_0=z_0$ for every $t\in G$.
\end{proof}

\section{Full answer}

The implemented action on $\HS(E)$ is
\[
 \mathcal T_t(B)=T_tBT_{t^{-1}}.
\]
The Hilbert--Schmidt space is a KH-module and $\mathcal T$ is an
$S$-covariant unitary action~\cite[Secs. 5--6]{EHK}.  Its fixed space is
$\HS_T(E)$, the invariant Hilbert--Schmidt operators.

\begin{theorem}\label{thm:main}
For every KH-dynamical system $(\A,S;E,T)$,
\[
 \AP_T(E)\subseteq E_{\rm ds}.
\]
Consequently, the following are equivalent:
\begin{enumerate}
\item $E$ has discrete spectrum, i.e. $E=E_{\rm ds}$;
\item $\AP_T(E)$ is order-dense in $E$.
\end{enumerate}
In particular, the open implication (ii)$\Rightarrow$(i) in Remark 3.8 of
\cite{HK} is true for arbitrary KH-dynamical systems.
\end{theorem}

\begin{proof}
Fix $x\in\AP_T(E)$ and $\varepsilon>0$.  Choose a finite set
$F\subseteq E$ satisfying~\eqref{eq:ap}, put $m=|F|$, and let
\[
 M=\ocl\operatorname{span}_{\A}F.
\]
This is a finite-rank KH-submodule of rank at most $m$~\cite[Sec. 2.5]{EHK}.
Let $P$ be its orthogonal projection.  The Pythagorean identity gives, for
$z\in E$ and $y\in M$,
\[
 |(I-P)z|^2\leq |z-y|^2.
\]
Taking the infimum over $y\in F$ in~\eqref{eq:ap} therefore yields
\begin{equation}\label{eq:Papprox}
 |(I-P)T_t x|\leq\varepsilon\one\qquad(t\in G).
\end{equation}

For $t\in G$ define the $\A$-linear orthogonal projection
\[
 Q_t=T_{t^{-1}}PT_t=\mathcal T_{t^{-1}}(P).
\]
Since $P$ has rank at most $m$, $P\in\HS(E)$ and
$|P|_{\HS}^2\leq m\one$.  Covariant unitarity gives the same bound for
every $Q_t$.  Moreover, by~\eqref{eq:Papprox},
\begin{equation}\label{eq:Qapprox}
 |x-Q_tx|
 =S_{t^{-1}}|(I-P)T_tx|
 \leq\varepsilon\one.
\end{equation}

Inside the lattice ball
$\{B\in\HS(E):|B|_{\HS}\leq\sqrt m\,\one\}$, let $C$ be the smallest
order-closed, convex, mix-closed set containing the orbit
$\{Q_t:t\in G\}$.  Equivalently, intersect all such subsets of the ball.
The family is nonempty because the ball itself is admissible.  The
implemented action permutes the orbit and preserves each of the defining
properties, so $\mathcal T_tC=C$.

Lemma~\ref{lem:fixed} supplies $A_\varepsilon\in C$ fixed by
$\mathcal T$, hence $A_\varepsilon\in\HS_T(E)$.  The constraint set
\[
 D_\varepsilon=\{B\in\HS(E):|x-Bx|\leq\varepsilon\one\}
\]
is order-closed, convex, and mix-closed.  For order-closedness, use the
bounded evaluation map
\[
 |Bx|\leq |B|_{\HS}|x|.
\]
By~\eqref{eq:Qapprox}, $D_\varepsilon$ contains every $Q_t$; hence
$C\subseteq D_\varepsilon$ and
\begin{equation}\label{eq:Aapprox}
 |x-A_\varepsilon x|\leq\varepsilon\one.
\end{equation}

The KH decomposition theorem identifies
\[
 E_{\rm ds}
 =\ocl\sum\{\ran B:B\in\HS_T(E)\}
 \qquad\text{\cite[Theorem 6.9]{EHK}}.
\]
Thus $A_\varepsilon x\in E_{\rm ds}$.  Apply~\eqref{eq:Aapprox} with
$\varepsilon=1/n$.  Then $A_{1/n}x\to x$ in norm (and hence in order),
so the order-closedness of $E_{\rm ds}$ gives $x\in E_{\rm ds}$.  This
proves $\AP_T(E)\subseteq E_{\rm ds}$.

If $\AP_T(E)$ is order-dense, the inclusion forces $E_{\rm ds}=E$.
Conversely, every vector in an invariant finite-rank KH-submodule has a
uniformly totally order-bounded orbit~\cite[Proposition 2.5]{HK}; the
definition of $E_{\rm ds}$ then gives order-density of $\AP_T(E)$ whenever
$E=E_{\rm ds}$.
\end{proof}

\section{Verification, scope, and novelty}

\begin{itemize}
\item The proof uses no global $L^2$ realization, invariant measure, group
amenability, or topology on $G$.
\item The fixed-point lemma was checked directly from order completeness,
mixing, convexity, and the parallelogram identity.  The critical evaluation
map is order-continuous because its lattice norm is dominated by
$|B|_{\HS}|x|$.
\item A tempting rank-one convexification can collapse to zero when supports
move under $S$.  The conjugate-projection construction avoids this: every
operator in the initial orbit approximately fixes the same $x$, and that
closed cyclic convex constraint survives the fixed-point step.
\item The conclusion is stronger than the requested density implication:
each conditionally almost-periodic vector lies in $E_{\rm ds}$.
\item Bounded novelty searches on 13 August 2026 covered the run indexes,
the revised arXiv source, and web/arXiv phrases combining
``KH-dynamical system'', ``conditionally almost periodic'', ``order-dense'',
``compact/discrete spectrum'', and ``cyclic convex fixed point''.  No later
answer to this exact abstract implication was found.  The source revision of
9 February 2026 still states it as open.  Novelty confidence is moderate:
the cyclic convex fixed-point lemma may be known in Banach--Kantorovich
language even though this application was not located.
\end{itemize}

\paragraph{Human-review recommendation.}
Prioritize the three interfaces in the proof: (1) the least lattice-norm
construction in Lemma~\ref{lem:fixed}; (2) invariance of the smallest bounded
order-closed convex mix-closed hull under the semilinear implemented action;
and (3) the use of the Edeko--Haase--Kreidler characterization of
$E_{\rm ds}$ by ranges of invariant Hilbert--Schmidt operators.  No other
unproved dependency remains.

\begin{thebibliography}{9}
\bibitem{HK}
M. Haase and H. Kreidler,
\emph{Precompactness notions in Kaplansky--Hilbert modules and extensions
with discrete spectrum}, arXiv:2505.14132v3, 2026.

\bibitem{EHK}
N. Edeko, M. Haase, and H. Kreidler,
\emph{A decomposition theorem for unitary group representations on
Kaplansky--Hilbert modules and the Furstenberg--Zimmer structure theorem},
Analysis Mathematica 50 (2024), 377--454; arXiv:2104.04865.
\end{thebibliography}

\end{document}
